% Source: Wald GR, physical PDF p. 153
%         (printed p. 140).
% Coverage: Figure 6.3, effective potential for timelike geodesics with \(L^2=6M^2\).
% Figures/tables/footnotes: Figure 6.3; no tables or footnotes.
% Status: source-reviewed and compile-clean.
% Uncertainties: none.
\begingroup
\renewcommand{\thefigure}{6.3}
\begin{figure}[htbp]
\centering
\begin{tikzpicture}[x=1.0cm,y=.78cm,>=latex,font=\small]
  \draw[->] (0,-3.25) -- (6.25,-3.25) node[right] {\(r\)};
  \draw[->] (0,-3.25) -- (0,1.35) node[above] {\(V\)};
  \foreach \x/\lab in {1.6/\(2M\),3.3/\(4M\),5.0/\(6M\)}
    \draw (\x,-3.18) -- (\x,-3.32) node[below=3pt] {\lab};
  \foreach \y/\lab in {-3/-3,-2/-2,-1/-1,0/0,1/1}
    \draw (-.07,\y) -- (.07,\y) node[left=3pt] {\lab};
  \draw[thick,smooth] plot coordinates {
    (.38,-3.22) (.48,-2.20) (.63,-1.08) (.84,-.36) (1.10,-.10)
    (1.60,-.01) (2.05,.18) (3.00,.31) (4.20,.37) (6.0,.38)};
\end{tikzpicture}
\caption{The effective potential, \(V\), for timelike geodesics in the case
\(L^2=6M^2\).}
\label{fig:6.3}
\end{figure}
\endgroup
